\chapter{Introduzione}
\label{ch:introduzione}

\section{Contesto generale sull'asset allocation e la gestione del rischio}

L'ottimizzazione del portafoglio e la gestione del rischio rappresentano due dei pilastri fondamentali della finanza quantitativa moderna. A partire dal lavoro pionieristico di Harry Markowitz del 1952 sulla \emph{Modern Portfolio Theory} (MPT), il problema dell'allocazione ottimale delle risorse è stato formalizzato come un problema di ottimizzazione vincolata a due obiettivi contrapposti: la massimizzazione del rendimento atteso e la contestuale minimizzazione della varianza del portafoglio. In questo paradigma economico, la matrice di covarianza e, in modo equivalente, la matrice di correlazione tra i rendimenti delle diverse attività finanziarie giocano un ruolo matematico centrale, agendo come l'operatore lineare che quantifica il rischio sistematico e i benefici della diversificazione.

Parallelamente, l'evoluzione dei mercati globali e l'introduzione di framework normativi sempre più stringenti hanno spinto l'industria finanziaria verso l'adozione di metriche di rischio più sofisticate della semplice varianza. Tra queste, il \emph{Value at Risk} (VaR) e l'\emph{Expected Shortfall} (ES) sono diventati gli standard \emph{de facto} per la quantificazione delle perdite potenziali in condizioni di mercato avverse entro un determinato orizzonte temporale e livello di confidenza. 

La stima accurata di tali metriche di rischio di coda si affida ampiamente a metodologie di simulazione, tra cui spiccano le simulazioni Monte Carlo. L'efficacia e l'affidabilità di questi motori di simulazione dipendono in maniera critica dalla qualità e dal realismo degli scenari di mercato generati. Se i modelli matematici sottostanti non riescono a catturare fedelmente le interdipendenze e la struttura correlativa degli asset finanziari, le stime del rischio risultano inevitabilmente distorte, conducendo a decisioni di \emph{asset allocation} sub-ottimali e a una potenziale sottostima delle minacce sistemiche.

\section{Definizione del problema di ricerca e motivazione dell'uso dei modelli generativi profondi}

Il principale ostacolo nell'applicazione pratica dei modelli classici risiede nella natura intrinseca delle matrici di correlazione empiriche. Quando stimata a partire da serie storiche campionarie di rendimenti, la matrice di correlazione soffre in modo significativo del problema del rumore statistico e della non-stazionarietà dei mercati. Quando il numero di asset analizzati ($N$) è dello stesso ordine di grandezza della lunghezza della serie storica ($T$), la matrice empirica diventa instabile, introducendo autovalori spuri che degradano le performance degli algoritmi di ottimizzazione, un fenomeno ampiamente studiato nell'ambito della \emph{Random Matrix Theory} (RMT).

Nel presente lavoro, l'analisi si concentra sui rendimenti giornalieri dei titoli componenti l'indice S\&P 500 estratti lungo un ampio orizzonte temporale ventennale, compreso tra il 2000 e il 2020. Per mitigare la complessità computazionale e l'instabilità numerica associate alla dimensionalità dell'intero indice, mantenendo al contempo una fedele rappresentazione macroeconomica, la ricerca isola un paniere target di 100 asset. Al fine di evitare distorsioni e preservare la struttura intrinseca del mercato, la selezione viene eseguita tramite un algoritmo di campionamento stratificato, il quale garantisce il rispetto delle percentuali originali dei settori definiti dal \emph{Global Industry Classification Standard} (GICS).

I metodi tradizionali per gestire il rumore in tali matrici (come lo \emph{shrinkage} o il \emph{bootstrapping} storico) presentano forti limiti: tendono a linearizzare le relazioni eliminando anomalie informative o rimangono vincolati alla limitatezza storica, impedendo l'esplorazione di scenari estremi mai osservati precedentemente ma statisticamente plausibili. In questo panorama, i modelli generativi profondi (\emph{Deep Generative Models}) offrono un cambio di paradigma radicale. Architetture basate sui \emph{Variational Autoencoders} (VAE) consentono di mappare lo spazio ad alta dimensionalità delle matrici di correlazione empiriche in uno spazio latente continuo e regolarizzato. 

Per ottimizzare l'efficienza computazionale ed eliminare le ridondanze geometriche dovute alla simmetria intrinseca e alla diagonale unitaria, l'input dei modelli viene ridotto alla sola componente triangolare inferiore linearizzata delle matrici estratte dal sottoinsieme stratificato. Attraverso questo vincolo rappresentazionale, i modelli apprendono a comprimere e successivamente riprodurre l'intera struttura informativa. 

Al fine di mappare con rigore scientifico le capacità di questi strumenti, il framework introduce un doppio binario di validazione, sia qualitativo che quantitativo, applicato uniformemente a tutte le tecniche testate --- Principal Component Analysis (PCA), Autoencoder Lineari (Linear AE), Autoencoder Profondi (AE) e Variational Autoencoder (VAE) --- sia in fase di ricostruzione storica che, nel caso del VAE, in fase di generazione stocastica di nuove matrici campionate dallo spazio latente. L'integrazione di metriche di errore matematico classiche combinate con metriche topologiche finanziarie permette di valutare non solo l'accuratezza puntuale dei modelli, ma anche la conservazione delle proprietà strutturali profonde del mercato.

\section{Obiettivi della tesi}

Il presente lavoro di tesi si propone di progettare, implementare e validare un framework basato su architetture neurali profonde per la compressione, ricostruzione e generazione stocastica di matrici di correlazione finanziaria realistiche, finalizzato al miglioramento delle simulazioni di rischio.

Gli obiettivi specifici della ricerca possono essere declinati nei seguenti punti:
\begin{enumerate}
    \item \textbf{Sviluppo della pipeline di pre-processing e campionamento:} Definire una procedura rigorosa per l'estrazione e la pulizia dei rendimenti dei titoli dell'S\&P 500 (2000-2020), implementando il campionamento stratificato basato sui settori GICS per isolare i 100 asset di studio ed effettuando la scomposizione per estrarre le sole matrici triangolari inferiori come input ottimizzato.
    \item \textbf{Studio comparativo e validazione multidimensionale:} Analizzare ed evidenziare quantitativamente le performance strutturali di PCA, Linear AE, AE e VAE nella fase di ricostruzione storica, calcolando e confrontando sistematicamente metriche matematiche di errore rigide quali il \emph{Mean Squared Error} (MSE), il \emph{Mean Absolute Error} (MAE) e la norma di Frobenius. Tale impianto quantitativo viene integrato da una valutazione qualitativa e topologica --- applicata sia alle ricostruzioni storiche sia alle nuove matrici sintetizzate dal VAE --- mediante algoritmi di \emph{clustering} gerarchico per mappare visivamente la preservazione dei cluster settoriali e l'estrazione del \emph{Minimum Spanning Tree} (MST) per misurare la fedeltà della connettività strutturale degli asset.
    \item \textbf{Generazione stocastica e validazione tramite Stylized Facts:} Configurare l'impianto probabilistico del VAE per sintetizzare nuove matrici di correlazione tramite campionamento dello spazio latente, estendendo ad esse l'intero protocollo qualitativo/quantitativo e validandole rigorosamente rispetto ai cinque \emph{stylized facts} fondamentali del mercato rispetto alla \emph{Random Matrix Theory} (RMT).
    \item \textbf{Integrazione e stima del rischio:} Applicare le matrici sintetiche generate e validate dal VAE all'interno di scenari Monte Carlo per calcolare il Value at Risk e l'Expected Shortfall di portafogli reali, dimostrando l'efficacia del framework nel catturare i rischi di coda in modo più robusto rispetto ai metodi storici convenzionali.
\end{enumerate}